% META: {"id": "TNSI-HISTOIRE-INFORMATIQUE-RE-C01", "chapitre": "TNSI-HISTOIRE-INFORMATIQUE", "type_objet": "remediation", "capacites": ["T-HIST-01A"], "capacites_codes": ["C1"], "mode_creation": "ex_nihilo", "sources_inspiration": [], "status": "generated"}

%---------------------------------------------------------------
% FICHE DE REMEDIATION — C01 : confondre l'idee et la machine
%---------------------------------------------------------------

\begin{exercice}{TNSI-HIST-RE-C01-EX1}{1}{10}
Un élève place, sur sa frise, la construction des premiers ordinateurs électroniques à
programme enregistré \textbf{avant} la formalisation de la machine universelle
par Alan Turing, au motif qu'« il faut bien une machine avant d'avoir une
théorie ».
\begin{enumerate}
  \item Situer chacun des deux événements dans le temps, à la décennie près.
  \item L'ordre proposé est-il le bon ? Citer un troisième repère, postérieur
    aux deux premiers, et le placer.
  \item Expliquer pourquoi une formalisation théorique peut précéder la machine
    qui la réalise, alors que l'intuition inverse paraît naturelle.
\end{enumerate}
\end{exercice}

% DOCUMENTARY-REVIEW: les dates relevent des documents, pas d'un tri de constantes.
% Turing, 1936, paragraphe 6 : https://www.cs.virginia.edu/~robins/Turing_Paper_1936.pdf
% EDVAC, 30 juin 1945, paragraphes 2.3-2.5 : https://real.mtak.hu/170042/1/Firstdraft.pdf
% Baby : https://curation.cs.manchester.ac.uk/computer50/www.computer50.org/mark1/new.baby.html

\begin{corrige}{TNSI-HIST-RE-C01-EX1}
\textbf{1.} La machine universelle est formalisée par Turing en $1936$, dans les
années trente ; les premiers ordinateurs électroniques à programme enregistré sont construits à
la fin des années quarante, autour de $1948$.

\textbf{2.} L'ordre proposé est \textbf{faux} : la formalisation précède les
machines d'une douzaine d'années. Un troisième repère commode est la mise en
service d'ARPANET, en $1969$, soit une vingtaine d'années après les premières
machines. La frise correcte est donc : Turing ($1936$), premiers ordinateurs électroniques à
programme enregistré ($1948$), ARPANET ($1969$).

\textbf{3.} Une théorie peut définir des objets et étudier leurs propriétés sans disposer d'un appareil qui les réalise. Turing décrit en 1936 un \textbf{modèle logique du calcul}, pour étudier ce qui peut être calculé par une procédure mécanique. La machine universelle de ce modèle peut simuler d'autres machines dont on lui fournit la description.

Cela ne signifie pas qu'il donne déjà tous les choix de construction d'un ordinateur. Le rapport EDVAC de \textbf{1945} décrit notamment une organisation où instructions et données sont stockées dans une même mémoire ; le Manchester Baby en \textbf{1948} réalise électroniquement le programme enregistré. Modèle théorique, choix d'architecture et réalisation matérielle sont donc des repères distincts. Leur chronologie ne prouve ni qu'une seule difficulté technique restait à résoudre, ni que toute l'informatique serait née d'une cause unique.
\end{corrige}
